% Isolated final-INCOMPLETE proposal B178. All prior proof/evidence bytes remain.
\subsection*{The paired contact and the full two-source relation kernel}
\statement{B.178}{Proved}{An actual paired-contact quotient retains both signed source actions and separates a prescribed word from the actual source kernel, strictly refining the joint comparison bound without deciding fixed-source surjectivity}
\label{stmt:B178}
Use the original filled group \(\Gamma _0\), its actual index-two
subgroup \(\Gamma _0^*\), and the marked source presentation
\[
 P_0=(\mathbb Z_A\times\mathbb Z_m)*F(B,D,E)
       \longrightarrow H_5
\]
of B.170--B.177. Let \(C_0^{\rm ord}\) be its ordinary kernel and
\(C_0\) the kernel of the actual compact epimorphism
\(\widehat P_0\to\mathcal H_5\). Commutators are
\([g,h]=ghg^{-1}h^{-1}\). There is an actual quotient of
\(\Gamma _0^*\) with the following full source-saturated presentation:
\begin{equation}\label{eq:B178-target}
\begin{gathered}
 F_5=F(a,U,k,q,y),\qquad S=F(m,V),\qquad K_0=F_5\rtimes S,\\
 c=U^{-1}aU,\qquad p=kak^{-1},\qquad t_U=kUk^{-1},
 \qquad H_D=\langle S,c,p\rangle,\\
 K_1=K_0*_{H_D}(H_D\times\mathbb Z_{D_0}),\\
 L_R=\langle S,a,U,p,D_0t_UD_0^{-1}\rangle\le K_1,\qquad
 G_{\rm pair}=K_1*_{L_R}(L_R\times\mathbb Z_R).
\end{gathered}
\end{equation}
The full signed action is printed below. The original contact images are
\begin{equation}\label{eq:B178-contacts}
 A_s\mapsto aR,\qquad B_s\mapsto a^{-1}U^{-1}a,\qquad
 D_s\mapsto UD_0U^{-1},\qquad E_s\mapsto1,
 \qquad m\mapsto m.
\end{equation}
The actual contact word
\[
                        w_2=[[m,D_s],[A_s,B_s]]
\]
survives in a finite quotient of the original \(\Gamma _0\).

In \(\widehat P_0\), define the closed normal subgroups
\[
 R_1=\overline{\langle\!\langle[m,D]\rangle\!\rangle},\qquad
 R_3=\overline{\langle\!\langle E,[A,B],[m,BDB^{-1}]
                                      \rangle\!\rangle}.
\]
Then B.170 and B.177's actual marked comparison maps satisfy
\begin{equation}\label{eq:B178-strict}
                 C_0\subsetneq R_1\cap R_3,
                 \qquad w_2\in(R_1\cap R_3)\setminus C_0.
\end{equation}
The corresponding ordinary inclusion is also strict, with the same word.
This is a statement about the internal relation kernel of the known
contact image. It neither identifies that kernel nor proves or refutes
surjectivity of the contact image onto the remaining fixed/cohomology fiber.

\begin{proof}
\emph{The actual deletion and all initial nodes.}
Use B.174's original free basis
\[
 F_{11}=F(a,B,d,e,f,j,\bar B,z,\bar e,g,x),
                  \qquad j=\bar a,\quad z=\bar d.
\]
The character \(\chi(f)=\chi(g)=1\in C_2\), zero on the other
letters, gives the actual even subgroup and free section
\(\langle m,v^*\rangle\), \(v^*=f^{-1}v\). B.174 proves that
\(H_5,N\) lie in this index-two subgroup and that all four signed
source operators preserve every positive family label with the stated
sheet shift. Retain the even \(a,B,j,z\) cores and the five coarse
Schreier generators; delete every other family core. This is an actual
source-invariant free rank-nine quotient of the even fiber. In
particular odd-core deletion is not imposed in the entire \(F_{11}\).

For the original numbered node list, the surviving whole initial rows
are exactly \(0,1,9,11,20,21,29,31\). All other initial words and all
forty odd-deck words vanish. Put
\[
 U=aB^{-1}a^{-1},\quad R=xjx^{-1},\quad D_0=xzx^{-1},
                       \quad c=U^{-1}aU.
\]
Here are both complete factors of those eight surviving commutators:
\[
\begin{array}{c|c|c}
 0&a&a^{-1}ja\\
 1&j&j^{-1}aj\\
 9&j&j^{-1}aBa^{-1}j\\
 11&j^{-1}aj&z\\
 20&c^{-1}&cR^{-1}c^{-1}\\
 21&j^{-1}&jx^{-1}c^{-1}xj^{-1}\\
 29&j^{-1}&jx^{-1}Uxj^{-1}\\
 31&jx^{-1}c^{-1}xj^{-1}&jz^{-1}j^{-1}
\end{array}
\]
Thus the complete initial relation set is equivalently
\begin{equation}\label{eq:B178-initial}
 [j,a]=[j,B]=[z,a]=[R,a]=[R,B]=[D_0,c]=1.
\end{equation}
Indeed the first four rows give the first three centralizations;
the next three give \([R,c]=[R,U]=1\), equivalent to
\([R,a]=[R,B]=1\); the last gives \([D_0,c]=1\).
Every signed source translate of these relations is retained below.
The full original contact words, before these centralizations, are
\[
 s_1=U,\quad s_2=1,\quad s_3=D_0^{-1},\quad
 s_4=UR^{-1}c^{-1}R(U^{-1}a)c^{-1}R^{-1}.
\]
The original marking is
\[
 A_s=s_1s_4^{-1}s_1^{-1},\quad
 B_s=A_s^{-1}s_1^{-1}A_s,\quad
 D_s=s_2s_1s_3^{-1}s_1^{-1}s_2^{-1},\quad E_s=s_2^{-1}.
\]
Substitution using only \([R,a]=[R,U]=1\) gives
\eqref{eq:B178-contacts}. No contact whisker has been discarded.

\emph{The complete raw second-source calculation.}
For this calculation first delete the whole families
\(d,e,\bar B,\bar e\), giving \(F(a,B,j,z,f,g,x)\), before any
even-sheet specialization. Set
\[
 Q=z^{-1}gz,\qquad P=xQ^{-1}x^{-1},\qquad L=U^{-1}fU.
\]
The full original rows, with their marked family conjugators, reduce to
\begin{equation}\label{eq:B178-raw}
\begin{gathered}
 h_a=PUP^{-1}U^{-1}f^{-1},\qquad
 h_B=h_a a^{-1}fUPa,\\
 h_j=L^{-1},\qquad h_z=L^{-1}z^{-1},\qquad
 v(x)=xQ^{-1}L,\\
 v(f)=PUP^{-1}U^{-1}P^{-1},\qquad
 v(g)=\operatorname{Ad}(L^{-1}zQx^{-1})(L).
\end{gathered}
\end{equation}
As usual \(v(a)=h_aah_a^{-1}\), and likewise for the other family
cores. A terminal own-core factor is removed only in its conjugation
of that same core. To make the whole raw-word reduction explicit, the
163-letter \(f\)-row is split at positions
\[
 0,21,29,50,71,92,113,134,142,163.
\]
Its nine blocks become
\(P,U,P^{-1},1,P^{-1},1,P,U^{-1},P^{-1}\).
The 239-letter \(g\)-row is split at
\[
 0,19,39,58,68,87,110,129,152,171,181,200,220,239.
\]
Its thirteen ordered blocks become
\[
\begin{gathered}
 L^{-1},\ Qx^{-1},\ L^{-1},\ 1,\ L,\
 xQ^{-1}zQx^{-1},\ L,\ xQ^{-1}z^{-1}Qx^{-1},\\
 L^{-1},\ 1,\ L,\ xQ^{-1},\ L.
\end{gathered}
\]
Their product is the last row of \eqref{eq:B178-raw}.
Substituting the complete \(h_B\) gives \(v(U)=PUP^{-1}\);
the same rows give \(v(R)=PRP^{-1}\) and \(v(D_0)=PD_0P^{-1}\).
The actual new even lift
\[
                  V=(f^{-1}P)^{-1}v^*=P^{-1}v
\]
is a change of section by an even fiber word. The source remains the
free group on \(m,V\), with the same source projection as \(m,v^*\).
Its action fixes \(U,R,D_0\), and, for \(G=jgj^{-1}\), has
\begin{equation}\label{eq:B178-normalized-raw}
\begin{gathered}
 V(f)=UP^{-1}U^{-1},\qquad
 V(a)=TaT^{-1},\quad T=UP^{-1}U^{-1}f^{-1},\\
 V(G)=\operatorname{Ad}(P^{-1}L^{-1}jzQx^{-1})(L).
\end{gathered}
\end{equation}
The lift \(V\) is not asserted to be a contact generator or a fixed
point of the original involution.

\emph{The actual invariant square quotient.}
Let \(F=afa^{-1}\), \(G=jgj^{-1}\). The normalized original \(m\)-rows are
\[
                     m(F)=F^{-1}GF,\qquad m(G)=F.
\]
Thus both signs preserve the normal subgroup generated by the even
words \(F^2,G^2\). Before imposing either square,
\[
                     V(F)=\operatorname{Ad}(V(a)Uxz^{-1}j^{-1})(G).
\]
This conjugator is even. Since \(L=U^{-1}a^{-1}FaU\), the other row
in \eqref{eq:B178-normalized-raw} makes \(V(G)\) a conjugate of \(F\).
Its initially odd conjugator
\[
                  C=(P^{-1}L^{-1}jzQx^{-1})U^{-1}a^{-1}
\]
can be replaced by the even \(CF\), with the same conjugate of \(F\).
The image of the square normal subgroup is therefore normal in the
even fiber, lies in that subgroup and contains its two generators up
to even conjugacy. It equals the subgroup, proving inverse invariance.

Use the reversible normalized even basis
\[
 a,U,R,D_0,t=F^2,u=GF^{-1},w_o=FG,x,y=FxF^{-1}.
\]
The invariant relations \(F^2=G^2=1\) are exactly \(t=1,uw_o=1\).
They leave a free rank-seven fiber. Put
\[
                 q=u,\qquad k=(RD_0)^{-1}x,\qquad H=Ukqy^{-1}.
\]
Its free basis is \(a,U,R,D_0,k,q,y\). These are relations of a further
actual quotient, not additional relations asserted in \(\Gamma _0\).

\emph{Both full signed actions on the retained free factor.}
The \(m\)-action fixes \(a,R,D_0\) and is
\begin{equation}\label{eq:B178-m-action}
 m(U)=q^{-1}Uq,\qquad m(k)=kq,\qquad m(q)=q,\qquad m(y)=q^{-1}y.
\end{equation}
The \(V\)-action fixes \(U,R,D_0\). Writing primes for its other
images, its complete ordered rows are
\begin{equation}\label{eq:B178-V-action}
\begin{aligned}
 a'&=HaH^{-1},\qquad &k'&=kU^{-1}a^{-1}H^{-1}U,\\
 y'&=Hay a(a')^{-1},&q'&=(k')^{-1}U^{-1}a'H(a')^{-1}y'.
\end{aligned}
\end{equation}
The primed letters on the right mean substitution of preceding rows.

Here is the even-sheet derivation retaining the additional \(a\)-factors.
After the two squares let \(u_o=gf^{-1}\), \(y_o=fxf^{-1}\).
A path on sheet zero or one rewrites \(f\) as the empty word with a
sheet change; \(g^{\pm1}\) contributes \(u_o\) on sheet zero and
\(u_o^{-1}\) on sheet one, again changing sheets. The letter \(x\)
contributes \(x\) or \(y_o\); a retained color contributes itself on
sheet zero and the empty word on sheet one. Negative letters use the
inverse path. The exact basis changes are
\[
 q=ju_oa^{-1},\qquad y=ay_oa^{-1},\qquad
 p_o=P|_0=kqy^{-1}a,\qquad P|_1=p_o^{-1}.
\]
The two \(L\)-words are \(U^{-1},U\). In particular
\(T|_0=Up_o=Ha\), not \(H\); its trailing \(a\) cancels only in
its conjugation of \(a\). The two \(V(F)\)-words are
\(Ha^2,a^{-2}H^{-1}\), and the two \(V(x)=xLP\)-words are
\(xU^{-1}p_o^{-1},y_oUp_o\). They give \(a',k',y'\) above.
The sheet-zero \(V(G)\)-word is
\[
                 p_oU(ju_o)y_o^{-1}Up_oy_oUp_o.
\]
Multiply it by the sheet-one \(V(F)^{-1}\)-word
\(a^{-2}H^{-1}\). Substituting the displayed changes gives
\[
 q'=p_oUqy^{-1}a(a')^{-1}y'
                      =(k')^{-1}U^{-1}a'H(a')^{-1}y'.
\]
This verifies all four rows. Setting \(a=1\) recovers B.175's rows,
but that specialization is not used to omit any factor.

The inverse is equally literal. From the output letters first form
\[
                   H=(a')^{-1}Uk'q'(y')^{-1}a',
\]
and recover, in the indicated order,
\begin{equation}\label{eq:B178-V-inverse}
 a=H^{-1}a'H,\qquad k=k'U^{-1}a'HU,\qquad
 y=a^{-1}H^{-1}y'a'a^{-1},\qquad q=k^{-1}U^{-1}Hy.
\end{equation}
The \(q'\)-row gives the first expression for \(H\); the next three
rows recover \(a,k,y\); the last enforces \(H=Ukqy^{-1}\).
Substitution in either direction gives the identity. The inverse of
\eqref{eq:B178-m-action} sends \(U,k,q,y\) to
\(qUq^{-1},kq^{-1},q,qy\), respectively. Thus both signed sources
preserve the whole free factor \(F_5\), and \(K_0=F_5\rtimes S\)
is the actual semidirect product asserted in the statement.

\emph{Both inverse whole-source presentations.}
The minus relations in \eqref{eq:B178-initial} say that \(R\)
centralizes \(xax^{-1},xUx^{-1}\), and \(D_0\) centralizes
\(xax^{-1}\). The other relations say that \(R\) centralizes
\(a,U\), and \(D_0\) centralizes \(c\).
Substitute \(x=RD_0k\). The first \(R\)-relation becomes
\([R,D_0pD_0^{-1}]=1\). Together with the \(D_0\)-relation it is
equivalent to \([D_0,p]=[R,p]=1\): conjugating \(D_0pD_0^{-1}\)
by \(R\) leaves it unchanged, so the second relation is its
centralization by \(D_0\), which implies \([D_0,p]=1\), and the
first then becomes \([R,p]=1\). The other \(R\)-relation becomes
\([R,D_0t_UD_0^{-1}]=1\). All implications reverse.

Both sources fix \(R,D_0\). Centralizing \(S\) and the displayed
finite lists therefore imposes precisely every signed source translate
of every initial node. This proves \eqref{eq:B178-target} in both
directions. The original source lift is recovered by
\(v^*=(f^{-1}P)V\), with its actual even fiber word. Every edge group
is the indicated actual subgroup of the preceding vertex; no freeness
or intrinsic completion of an infinite edge group is asserted.
Ordinary amalgam normal form embeds the successive bases. Sending
\(R\) to one gives a retraction \(G_{\rm pair}\to K_1\).

\emph{An invariant of the entire signed edge orbit.}
Kill \(a\) in \(F_5\), obtaining \(T_0=F(U,k,q,y)\).
The ordinary kernel \(N_a\) has the literal free basis
\[
                         a_g=gag^{-1},\qquad g\in T_0.
\]
Write \(m_0,V_0\) for the quotient actions. Then
\[
 m(a_g)=a_{m_0(g)},\qquad
 V(a_g)=n_g a_{V_0(g)H}n_g^{-1},\qquad n_g\in N_a,
\]
where \(V(g)=n_gV_0(g)\). The inverse has the inverse affine label
map. The corrections \(n_g\) are retained throughout.
Define integer characters on \(T_0\) by
\[
 \lambda _1(g)=\operatorname{exp}_U(g)+\operatorname{exp}_y(g),
 \qquad
 \lambda _2(g)=\operatorname{exp}_k(g)+\operatorname{exp}_y(g).
\]
On the ordered coordinates \((U,k,q,y)\), the \(V_0\)-columns are
\[
 (1,0,0,0),\quad(-1,0,-1,1),\quad
 (2,2,3,-2),\quad(1,1,1,0).
\]
The \(m_0\)-map sends \((u,k,q,y)\) to \((u,k,q+k-y,y)\).
Both characters are invariant under both signed actions, and
\(\lambda(H)=(0,0)\). Thus the full signed affine source action
preserves each label's pair of integer values, with no period cutoff.

The fiber of \(H_D=\langle S,c,p\rangle\) is generated by the
complete \(S\)-orbits of \(c=a_{U^{-1}}\) and \(p=a_k\).
It therefore lies in the \(N_a\)-normal closure of basis letters
with \(\lambda=(-1,0)\) or \((0,1)\). This assertion includes
all correction conjugators \(n_g\) and all inverse source words.
It is not a permutation-only model of the action.

\emph{The two nonedge tests and the full word.}
Put \(U_1=m(U)=q^{-1}Uq\) and \(h=U_1^{-1}U\). The two inner
contact commutators are
\[
 c_m=[m,D_s]=U_1D_0hD_0^{-1}U^{-1},\qquad
 s_m=[A_s,B_s]=[a,a^{-1}U^{-1}a]=U^{-1}a^{-1}Ua.
\]
The second identity uses \([R,a]=[R,U]=1\), so \(R\) centralizes
\(B_s\). Killing \(a\) sends \(H_D\) to the source section \(S\),
whereas it sends \(h\) to the nontrivial word
\(q^{-1}U^{-1}qU\) with trivial source component. Hence \(h\notin H_D\).
Also
\[
                       U^{-1}s_mU=a_{U^{-2}}^{-1}a_{U^{-1}}.
\]
The first letter has label \((-2,0)\), the second \((-1,0)\).
Killing all basis letters with either of the two edge values kills the
whole \(H_D\)-fiber and leaves the first letter nontrivial. This proves
\(U^{-1}s_mU\notin H_D\).
The full mixed word in \(K_1\) is
\begin{equation}\label{eq:B178-Britton}
 [c_m,s_m]
 =U_1D_0hD_0^{-1}(U^{-1}s_mU)
       D_0h^{-1}D_0^{-1}U_1^{-1}s_m^{-1}.
\end{equation}
The three possible pinches have coefficients
\(h,U^{-1}s_mU,h^{-1}\), all outside \(H_D\). This is a reduced
four-stable-letter word, nontrivial in \(K_1\) and hence in the full
source-saturated target. No one-source partial quotient is substituted.

\emph{Actual finite ambient separation.}
Both nonedge tests have simultaneous finite witnesses on the free fiber.
Let \(A_2=(C_2)^2\) and map the head letters by \(\lambda\bmod2\):
\[
 U\mapsto(1,0),\quad k\mapsto(0,1),\quad q\mapsto0,
                         \quad y\mapsto(1,1).
\]
Map \(a\) to the zero-index lamp in
\(W_2=\mathbb F_2[A_2]\rtimes A_2\).
Then \(a_g\) maps to the lamp at \(\lambda(g)\). Every \(N_a\)
correction maps into the abelian lamp subgroup. The entire edge fiber
therefore maps into the span of the lamps at \((1,0),(0,1)\), while
\(U^{-1}s_mU\) maps to the sum of the lamps at \((0,0),(1,0)\),
which lies outside that span. The edge fiber is not claimed to be killed.
For \(h\), the prescribed map
\[
            U\mapsto(12),\quad q\mapsto(23),\quad a,k,y\mapsto1
\]
to \(S_3\), with rightmost permutations first, gives \(h\mapsto(123)\).
It kills every conjugate of \(a\).

These are maps of \(F_5\) only. In fact \(k^2\) maps to one in
\(W_2\), while \(V(k^2)\) maps to the sum of the lamps at
\((1,1),(1,0)\), which is not one. The source action does not descend
to this small lamp group. Intersect the two fiber kernels and choose
a characteristic finite-index subgroup \(B_0\) of \(F_5\) inside
the intersection. Such refinements are cofinal: intersect the finitely
many subgroups of index at most the index of the given intersection.
The whole signed source action preserves \(B_0\) and has finite image
in \(\operatorname{Aut}(F_5/B_0)\). Thus there is an actual finite head
\[
 Q_0=(F_5/B_0)\rtimes
                \operatorname{im}\bigl(S\to\operatorname{Aut}(F_5/B_0)\bigr).
\]
Let \(H_0\) be the image of \(H_D\). Its elements with trivial
source component have fiber in the subgroup generated by the entire
source-orbit images of \(c,p\), because the finite source image has
its literal zero-fiber section. The two fiber witnesses therefore prove
that the images of \(h\) and \(U^{-1}s_mU\) both lie outside \(H_0\).
This is kernel refinement, not extension into the same small codomain.

Map \(K_1\) into the finite-vertex amalgam
\[
                     Z=Q_0*_{H_0}(H_0\times C_{2,D_0}).
\]
The word \eqref{eq:B178-Britton} remains reduced and nontrivial.
The map \(Z\to Q_0\times C_2\) is injective on both finite vertices,
so its kernel is a finite-index finitely generated free group, by its
free action on the ordinary amalgam tree. An element outside the kernel
is separated by this finite map. A nontrivial element inside is separated
by a finite quotient of the free kernel followed by a characteristic
finite-index refinement, normal in the whole amalgam. Thus \(Z\) is
residually finite and some finite quotient separates the specified word.

Compose \(G_{\rm pair}\to K_1\), \(R\mapsto1\), with these maps.
Pull the finite kernel back to \(\Gamma _0^*\), then take its finite
index normal core in the whole original \(\Gamma _0\). Its intersection
with \(\Gamma _0^*\) is contained in the separating kernel, so the
original contact word \(w_2\) still survives. This gives the asserted
finite ambient quotient. It uses no full-topology assertion for \(H_D\),
\(L_R\), or an infinite-index contact subgroup.

\emph{The precise compact conclusion and its limit.}
B.170 gives \(C_0\subseteq R_1\), and B.177 gives the actual compact
comparison with kernel \(R_3\). The same statements hold ordinarily.
The word \(w_2\) lies in \(R_1\) because its first inner commutator
does, and in \(R_3\) because its second inner commutator does.
Finite ambient survival puts it outside \(C_0\) and outside its ordinary
analogue. This proves \eqref{eq:B178-strict} and its ordinary version.

In B.164's notation, this word is already in \(\mathcal I_5\subseteq
\mathcal I\), so \(\partial(w_2)\) is a known-source class. Neither
its nontriviality nor the exact presentation of the comparison target
computes arbitrary classes in the whole kernel
\(H^1(\Delta,\mathcal N)\to H^1(\Delta,\mathcal K_a)\).
After taking a \(\Delta\)-stable finite core, a finite filled quotient
still kills \(\mathcal N\); at that generic finite coordinate the node
cohomology kernel is the singleton, even when the filled word survives.
Intersecting the separating kernel with its \(\Delta\)-image preserves
both the node-kernel containment and the finite separation. Thus this separator does not establish B.164's
cofinal surjectivity criterion or the separate original B.75 criterion.
\end{proof}

\begingroup\footnotesize\raggedright
The complete prior proof and its whole-source independent reviews are
bound in \path{work_DE/paired_A_full_source_mixed_survival38.md} and
\path{work_DE/paired_A_full_source_review_binding38.json}.
The sole prescribed source checker is
\path{work_CE/code/check_paired_A_full_source38.py}; its actual exit zero
and complete output are bound in
\path{work_CE/paired_A_full_source38_observed.json}.
Its 588 exact checks cover the original signed rows, all twenty-one
even-generator diagrams for four signs, both normalized basis inverses,
raw factor blocks, square conjugators, all forty nodes and decks, contacts,
integer invariants and the two specified finite-fiber witnesses.
The complete certificate is
\path{work_CE/next_symmetric_nodes/paired_A_full_source38.json}.
It explicitly records the failure of source descent to the small lamp
group. No numerical characteristic ambient quotient or final finite
amalgam quotient was constructed; their existence is proved above.
No further census, source relation, fixed-generation result or Main B
decision is inferred.
\par\endgroup
